% !TeX root = ../thesis.tex
\clearpage
\section{Query-Action Contract}
\label{app:query-action-contract}

This appendix shows the input and feedback exchanged through
\texttt{executeAthenaQuery}. The function definition is shared by the Nova Lite
and Haiku templates.

\subsection{Input}

Listing~\ref{lst:query-action-input} reproduces the
\texttt{executeAthenaQuery} function definition from the CloudFormation template.
The only parameter supplied by the agent is
\texttt{query}; database, workgroup and result-location settings are supplied
through the Lambda configuration.
\texttt{RequireConfirmation: DISABLED} allows the function to be invoked
without asking the user to approve each call. This setting is separate from
asking for clarification of an ambiguous request.

\begingroup
\lstset{language={},basicstyle=\footnotesize\ttfamily,breaklines=true,
  columns=fullflexible,keepspaces=true,showstringspaces=false,
  frame=single,captionpos=b,numbers=none}
\begin{minipage}{\textwidth}
\begin{lstlisting}[caption={Query-action definition.},label={lst:query-action-input}]
Name: executeAthenaQuery
Description: Execute one read-only Athena SELECT or WITH query.
RequireConfirmation: DISABLED
Parameters:
    query:
        Type: string
        Required: true
        Description: A single read-only Athena SQL SELECT or WITH query.
\end{lstlisting}
\end{minipage}

\subsection{Successful Response}

Listing~\ref{lst:query-action-success} shows the saved response to counting
artists in the catalogue (CB200-E001), from Nova's first repetition.
The query returned one row containing the count of 275 artists.

\begin{minipage}{\textwidth}
\begin{lstlisting}[caption={Saved artist-count response.},label={lst:query-action-success}]
{
  "queryExecutionId": "e6f07dca-dc19-4cca-b806-934f300c9d1e",
  "query": "SELECT COUNT(*) AS artist_count FROM artists",
  "columns": [
    "artist_count"
  ],
  "rows": [
    {
      "artist_count": "275"
    }
  ],
  "previewRowCount": 1,
  "previewTruncated": false
}
\end{lstlisting}
\end{minipage}

\texttt{previewRowCount} counts the returned preview rows.
\texttt{previewTruncated} indicates whether more rows exist beyond that preview.
The configured preview limit was 200 rows.
The execution identifier allowed the runner to retrieve the complete result
for offline evaluation.

\subsection{Error Response}

Validation and execution exceptions were caught by the handler, which returned
an \texttt{error} field with \texttt{responseState} set to \texttt{REPROMPT}.
This state passes the error message back to the model so that it can respond
or revise the query.

Listing~\ref{lst:query-action-error} reproduces the saved error payload for
listing invoices issued in January 2010 (CB200-E030), from Haiku's first
repetition. Its first query compared the timestamp column
\texttt{invoice\_date} with string literals such as \texttt{'2010-01-01'}.
Athena rejected the comparison because the operand types differed.

\begin{minipage}{\textwidth}
\begin{lstlisting}[caption={Saved date-comparison error.},label={lst:query-action-error}]
{
  "error": "Athena query failed: TYPE_MISMATCH: line 3:20: Cannot apply operator: timestamp(3) <= varchar(10)"
}
\end{lstlisting}
\end{minipage}
\endgroup

Haiku's next query cast \texttt{invoice\_date} to \texttt{DATE} and used
\texttt{DATE} literals for the range boundaries. It executed successfully
and returned seven invoice rows. This example shows the database diagnostic
available to the agent before it revised the SQL.
